%\chapter{Projekt Requirements}
%\section{Projekt Requirements}
%\label{sec:Projekt Requirements}
%
%\textcolor{red}{ICH WÜRDE DIE NACH FUNKTIONEN WIE NOTFALLBREMSE DISPLAY KI USW AUFTEILEN SO MACHT MAN AUCH BEI ADAS BEI STELLANTIS Jeder kann die Requirements seiner Aufgabe schireben} 
%
%\begin{longtable}{|p{1.7cm}|p{4.2cm}|p{8.8cm}|p{1.9cm}|}
%
%
%	\caption{Übersicht der Anforderungen für das Projekt}
%	\label{tab:anforderungen_ueberarbeitet} \\
%	\hline
%	\textbf{ID} & \textbf{Unterkategorie} & \textbf{Anforderungsbeschreibung} & \textbf{Priorität} \\
%	\hline
%	\endfirsthead
%	
%	\hline
%	\textbf{ID} & \textbf{Unterkategorie} & \textbf{Anforderungsbeschreibung} & \textbf{Priorität} \\
%	\hline
%	\endhead
%	
%	\hline
%	\multicolumn{4}{r}{Fortsetzung auf nächster Seite} \\
%	\endfoot
%	
%	\hline
%	\endlastfoot
%	
%	% =========================================================
%	% FUNKTIONALE ANFORDERUNGEN
%	% =========================================================
%	\multicolumn{4}{|l|}{\textbf{Funktionale Anforderungen}} \\
%	\hline
%	REQ-001 & Bewegung & Der Roboter muss vorwärts und rückwärts fahren können. & High \\
%	\hline
%	
%	% =========================================================
%	% NICHT-FUNKTIONALE ANFORDERUNGEN
%	% =========================================================
%	\multicolumn{4}{|l|}{\textbf{Nicht-funktionale Anforderungen (Performance/Technik)}} \\
%	\hline
%	REQ-101 & Reaktionszeit & Die Motorsteuerung muss eine Verzögerung von weniger als 0{,}5\,s zwischen Steuersignal und Reaktion aufweisen. & High \\
%	\hline
%
%	
%	% =========================================================
%	% SAFETY
%	% =========================================================
%	\multicolumn{4}{|l|}{\textbf{Sicherheitsanforderungen (Safety)}} \\
%	\hline
%	REQ-201 & Fail-Safe & Wird kein Fahrsignal erkannt, muss der Roboter selbstständig anhalten. & High \\
%	\hline
%
%	% =========================================================
%	% QUALITÄT
%	% =========================================================
%	\multicolumn{4}{|l|}{\textbf{Qualitätsanforderungen}} \\
%	\hline
%	REQ-301 & Tragstruktur & Die Struktur muss das gesamte Systemgewicht sicher tragen können. & High \\
%	\hline
%
%
%	% =========================================================
%	% WARTBARKEIT
%	% =========================================================
%	\multicolumn{4}{|l|}{\textbf{Wartbarkeit und Betrieb}} \\
%	\hline
%	REQ-401 & Bedienbarkeit & Die Bedienung muss einfach und selbsterklärend sein. & Medium \\
%
%	
%	% =========================================================
%	% DESIGN
%	% =========================================================
%	\multicolumn{4}{|l|}{\textbf{Design und Außenwirkung}} \\
%	\hline
%	REQ-501 & Visuelle Gestaltung & Das Erscheinungsbild soll der DHBW-Designsprache entsprechen. & Low \\
%	\hline
%
%	
%\end{longtable}
